\documentclass[11pt]{article}
\usepackage[margin=1in]{geometry}
% Common macros for the anonymity series (shared across papers).
% Keep this file small: each paper should still compile as a standalone artifact.

\usepackage{amsmath,amssymb,amsthm}
\usepackage{mathtools}
\usepackage{microtype}
\usepackage{enumitem}
\usepackage{booktabs}
\usepackage{hyperref}
\usepackage[nameinlink,noabbrev]{cleveref}

\newcommand{\RR}{\mathbb{R}}
\newcommand{\NN}{\mathbb{N}}
\newcommand{\EE}{\mathbb{E}}
\newcommand{\PP}{\mathbb{P}}

\DeclareMathOperator{\TV}{TV}
\DeclareMathOperator{\KL}{D_{\mathrm{KL}}}
\DeclareMathOperator{\MI}{I}

\newcommand{\one}{\mathbf{1}}

\newcommand{\defeq}{\stackrel{\mathrm{def}}{=}}

\newtheorem{theorem}{Theorem}
\newtheorem{proposition}{Proposition}
\newtheorem{lemma}{Lemma}
\newtheorem{corollary}{Corollary}
\theoremstyle{definition}
\newtheorem{definition}{Definition}
\theoremstyle{remark}
\newtheorem{remark}{Remark}

% A small "callout" macro for short papers.
\newcommand{\takeaway}[1]{\paragraph{Takeaway.}\emph{#1}}


\title{Posterior-Mass Inflation Anonymity:\\ Pointwise Maximal Leakage, Composition, and the True-Odds Boundary}
\author{}
\date{June 18, 2026 (archive rev0899)}

\begin{document}
\maketitle

\begin{abstract}
This note repairs a semantic error in the archive's endpoint vocabulary.
The quantity
$\log_2\!\bigl(\pi(i\mid o)/\pi(i)\bigr)=\log_2\!\bigl(P_i(o)/P(o)\bigr)$
is posterior-\emph{mass} inflation and, after maximizing over identities, is exactly pointwise maximal leakage (PML); it is not posterior-odds inflation.
We retain the valid pointwise, post-processing, composition, identification, and shortlist consequences of that quantity, derive its exact conversion to true one-vs-complement posterior odds when a prior cap prevents posterior saturation, and give a deterministic-reveal counterexample showing that no prior-independent finite odds conversion exists.
We also separate hockey-stick slack from a bad-transcript probability: a tail certificate implies a hockey-stick envelope, but the converse can fail by a factor of 51 even in a two-identity example.
Finally, we tighten the common-reference TV conversion by using the sharp singleton consequence $|P(o)-Q(o)|\le \TV(P,Q)$ rather than paying an unnecessary factor of two.
The result is a corrected endpoint theorem, not a deployment certificate.
\end{abstract}
\paragraph{Scope.}
We define a mechanism-agnostic endpoint for anonymity: a bound on realized posterior-mass inflation for every observable transcript, together with explicit rules for when that endpoint may or may not be translated into true odds.

\paragraph{Units.} Throughout, information budgets are measured in bits: we write $\log_2$ and $2^{\epsilon}$ for posterior-mass-inflation / PML caps. Imported results stated in nats are converted by dividing by $\ln 2$.
A release card must not mix units silently: every imported moment accountant, hockey-stick profile, or privacy-filter row declares its log base and converts to this paper's bit-valued endpoint before applying the formulas below.

\paragraph{Imports.}
Mechanism-specific scheduling, spectral, stop-time, and equalization contracts are black-box prerequisites rather than claims of this paper.\cite{series:schedcompiler,series:spectral,series:stoptimeaddendum,series:knapsacktransport}
The bounded transcript projection and repetition horizon remain explicit claim fields: \texttt{exposure\_nf\_id} (ENF-ID) and \texttt{tw\_id} (TW-ID).\cite{series:traceifaces,series:threatwindows}
Notation conventions follow the series notation sheet.\cite{series:notation}

\paragraph{Exports.}
This paper exports (i) posterior-mass inflation / PML as an endpoint metric, (ii) exact and hockey-stick sequential-composition bounds, (iii) operational identification and shortlist consequences, (iv) a guarded conversion to true one-vs-complement odds, and (v) corrected TV and R\'enyi conversion lemmas.
These outputs may feed Boss Fight accounting and certificate workflows only after their observation, prior, horizon, and evidence fields are pinned.\cite{series:endpointbridge}

\paragraph{Boundary.}
This note owns endpoint semantics and theorem conversions only.
It does not certify an anonymous-DHT deployment, select a prior, establish a transcript-equalization witness, size a repeated-use horizon, or serialize a public claim; those duties remain with Anonymity~B, Boss Fight~A, Anonymity~C, and Release~A.\cite{series:bossfightA,series:profiling,series:abomoinl,series:releaseA}
A downstream card that still labels $\log_2(P_i/P)$ as ``posterior odds'' is semantically stale even when its numerical posterior-mass or identification arithmetic happens to remain valid.

\paragraph{Later-terse-paper citation posture.}
Cite this note when the live issue is the PML/posterior-mass endpoint, the hockey-stick checkpoint, the true-odds conversion boundary, or a generic TV / moment conversion.
Cite Boss Fight~A for repeated-use sizing, Anonymity~B for anonymous-DHT profiling, Anonymity~D for an optional delegation witness-to-budget translation, and Anonymity~C plus Release~A for claim packaging and lineage.\cite{series:bossfightA,series:profiling,series:spectral2budget,series:abomoinl,series:releaseA}

\paragraph{Artifact policy.}
The theorem is source-complete. Archive-level compile receipts and hostile-review vectors are evidence about this source, not hidden premises of the mathematical claim.\cite{series:artifactpolicy}

\section{Introduction}
Anonymous systems often come with a stack of mechanism-specific arguments: timing padding, cover traffic, routing randomization, batching, and other defenses.
A certification layer still needs a compact endpoint that says what any such stack has proved about an observer's posterior.
The archive previously used the word ``odds'' for the ratio $\pi(i\mid o)/\pi(i)$.
That ratio compares posterior and prior \emph{probabilities}; true posterior odds compare $i$ against its complement and use $\pi(i\mid o)/(1-\pi(i\mid o))$.
Confusing the two is harmless only in a rare-secret approximation and can be catastrophic when the posterior approaches one.

This repair keeps the valid endpoint: a pointwise cap on posterior-mass inflation, exactly PML under the declared prior.
It then makes the true-odds conversion explicit, including its prior-dependent denominator, and separates hockey-stick divergence from tail probability.
Mechanism details remain outside this paper; the point is to make the exported theorem impossible to overread.

\section{Model and definitions}
Let $I\in\{1,\dots,N\}$ be a hidden identity with prior $\pi$, with $\pi(i)>0$ for every identity in scope, and let $O$ be any public observation (single output or full transcript).
Write $P_i$ for the law of $O$ given $I=i$, and $P=\sum_i \pi(i)P_i$ for the mixture (the public marginal).
We use $\TV(P,Q)=\tfrac12\|P-Q\|_1$ for total variation on finite alphabets.

\begin{definition}[Per-identity posterior-mass inflation and PML]\label{def:mass_inflation}
For any $o$ with $P(o)>0$, define
\[
m_i(o)\ \defeq\ \frac{\pi(i\mid o)}{\pi(i)}\ =\ \frac{P_i(o)}{P(o)},
\qquad
L_i(o)\ \defeq\ \log_2 m_i(o).
\]
The realized worst-case posterior-mass inflation is
\[
\ell(o)\ \defeq\ \max_i L_i(o).
\]
Under the same prior, $\ell(o)$ is exactly bit-valued pointwise maximal leakage at outcome $o$.
\end{definition}

\begin{definition}[Exact PML cap and hockey-stick checkpoint]\label{def:epsinf}
A mechanism has an \emph{$\varepsilon_\infty$ posterior-mass / PML cap} relative to $\pi$ if $\ell(o)\le \varepsilon$ for every $o$ with $P(o)>0$.
Separately, it has an \emph{$(\varepsilon,\delta)_{\mathrm{HS}}$ identity-vs-mixture hockey-stick envelope} if, for every event $A$ and every identity $i$,
\begin{equation}\label{eq:hockeystick}
P_i(A)\ \le\ 2^\varepsilon P(A)+\delta.
\end{equation}
Equivalently, $E_{2^\varepsilon}(P_i\|P)\le\delta$ for every $i$.
\end{definition}

\begin{remark}[What $\delta$ does and does not mean]\label{rem:delta-not-tail}
The hockey-stick slack is the weighted excess
\[
E_{2^\varepsilon}(P_i\|P)
=\sum_o\bigl(P_i(o)-2^\varepsilon P(o)\bigr)_+,
\]
not, in general, the probability that $L_i(O)>\varepsilon$.
It supports event inequalities, post-processing, and the checkpoint composition below, but it must not be advertised as a bad-transcript probability unless a separate tail theorem supplies that interpretation.\cite{LiuCuffVerdu2017,DworkRoth2014}
\end{remark}

\begin{proposition}[A tail certificate implies hockey-stick; the converse fails]\label{prop:tail-vs-hs}
If $P_i(L_i(O)>\varepsilon)\le\eta$, then $E_{2^\varepsilon}(P_i\|P)\le\eta$.
The converse is false, even for two identities and $\varepsilon=0$.
\end{proposition}

\begin{proof}
Let $H=\{o:L_i(o)>\varepsilon\}$.
For every event $A$,
$P_i(A\cap H^c)\le 2^\varepsilon P(A\cap H^c)$, while $P_i(A\cap H)\le P_i(H)\le\eta$; adding gives the hockey-stick inequality.
For the converse, take a uniform prior on two identities and outcomes $a,b$ with
$P_1=(0.51,0.49)$ and $P_2=(0.49,0.51)$.
Then $P=(0.5,0.5)$ and $E_1(P_i\|P)=\TV(P_i,P)=0.01$ for both identities, so the mechanism satisfies $(0,0.01)_{\mathrm{HS}}$.
But under $P_1$, $L_1(a)=\log_2(1.02)>0$ and $P_1(a)=0.51$; symmetrically for identity $2$.
Thus the positive-loss tail probability is $0.51$, not $0.01$.
\end{proof}

\section{Priors and contexts as certificate inputs}\label{sec:priors}
Posterior-mass inflation is defined relative to a prior $\pi$ because the public marginal $P=\sum_i \pi(i)P_i$ appears in the Bayes factor $P_i/P$.
This is a feature: a deployment's \emph{context} (who could plausibly be the requester) is part of the privacy surface.
Accordingly, in our series we treat priors and eligibility constraints as first-class certificate inputs (Anonymity~C) \cite{series:abomoinl}.

\paragraph{Certification practice.}
A paper should state which prior or prior family its claimed budget targets.
Two common, conservative patterns are:
\begin{enumerate}[leftmargin=*]
\item \textbf{Uniform over an eligibility pool.} Declare an eligibility interface (e.g.\ ``clients active in this epoch''), and take $\pi$ uniform over that pool.
Then the operational corollaries in \Cref{prop:shortlist,prop:property} translate directly into lower bounds on the number of plausible requesters.
\item \textbf{Capped prior mass.} Declare a cap $\pi(i)\le \pi_{\max}$ for all identities in scope.
Under an $\varepsilon_\infty$ posterior-mass/PML cap, every posterior mass satisfies $\pi(i\mid o)\le 2^{\varepsilon}\pi_{\max}$, so deployments can report a policy-friendly statement such as ``no single identity exceeds $p^\star$ posterior mass'' by enforcing $2^{\varepsilon}\pi_{\max}\le p^\star$.
\end{enumerate}
In both cases, \emph{eligibility metadata must be treated as a public interface} and budgeted (or hidden) rather than silently absorbed into the prior.

\begin{proposition}[Eligibility/selection disclosures add a prior-shift term]\label{prop:selection}
Let $C\subseteq\{1,\dots,N\}$ with $\pi(C)>0$ and let $E\defeq\one\{I\in C\}$ denote an eligibility/selection bit that becomes public (e.g.\ an observer learns which shard/committee the requester must belong to).
Let $P_C(o)\defeq\sum_{j\in C}\pi(j\mid C)P_j(o)$ be the public marginal conditioned on $I\in C$.
Then for any $i\in C$ and any observation $o$ with $P_C(o)>0$, the posterior-mass inflation for the combined disclosure $(E{=}1,O{=}o)$ satisfies
\[
\log_2\frac{\pi(i\mid E{=}1,o)}{\pi(i)}\ =\ -\log_2 \pi(C)\ +\ \log_2\frac{P_i(o)}{P_C(o)}.
\]
In particular, if the mechanism has an $\varepsilon_\infty$ posterior-mass/PML cap relative to the conditional prior $\pi(\cdot\mid C)$ (i.e., $\max_{i\in C}\log_2\tfrac{P_i(o)}{P_C(o)}\le\varepsilon$ for all $o$), then relative to the original prior it has a posterior-mass/PML cap of $\varepsilon+\log_2\tfrac{1}{\pi(C)}$ for identities in $C$ (equivalently, $\varepsilon-\log_2\pi(C)$).
Under a uniform prior this adds the familiar term $\log_2(N/|C|)$.
\end{proposition}

\begin{proof}
For $i\in C$, $P_i(E{=}1)=1$ while $P(E{=}1)=\pi(C)$.
Moreover $P(E{=}1,O=o)=\pi(C)\,P_C(o)$.
Therefore
$\log_2\tfrac{P_i(E{=}1,O=o)}{P(E{=}1,O=o)}=\log_2\tfrac{P_i(o)}{\pi(C)P_C(o)}=-\log_2\pi(C)+\log_2\tfrac{P_i(o)}{P_C(o)}$.
Finally, $\pi(i\mid E{=}1,o)/\pi(i)=P_i(E{=}1,O=o)/P(E{=}1,O=o)$ by Bayes.
\end{proof}

\begin{remark}[Context-aware leakage bounds]
Recent pointwise-maximal-leakage analyses show that even for fixed mechanisms, leakage can be substantially smaller under a committed prior model (e.g., minimum-probability assumptions) than under worst-case, context-free bounds \cite{ZhaoSaeidianOechtering2026ContextAware}.
Posterior-mass/PML anonymity makes this dependence explicit by treating the prior (or prior family) as a certificate input.
\end{remark}

\section{Posterior mass, PML, and true posterior odds}\label{sec:true-odds}
The PML quantity $m_i(o)=P_i(o)/P(o)$ is a posterior-to-prior \emph{probability} factor.
True one-vs-complement odds use a different denominator.
For $p_i=\pi(i)<1$, define the complement channel
\[
P_{-i}(o)\ \defeq\ \sum_{j\ne i}\frac{\pi(j)}{1-p_i}P_j(o)
\]
and the true odds-inflation factor
\[
R_i(o)\ \defeq\
\frac{\pi(i\mid o)/(1-\pi(i\mid o))}{\pi(i)/(1-\pi(i))}
=\frac{P_i(o)}{P_{-i}(o)},
\]
with $R_i(o)=+\infty$ when the posterior of $i$ is one and the prior is below one.

\begin{proposition}[Exact conversion from posterior mass to true odds]\label{prop:mass-to-odds}
Let $p=\pi(i)$ and $m=m_i(o)$.
If $pm<1$, then
\[
R_i(o)=m\,\frac{1-p}{1-pm}.
\]
Consequently, if $\ell(o)\le\varepsilon$, $\pi(i)\le p_{\max}$ for every identity, and $2^\varepsilon p_{\max}<1$, then
\[
\log_2 R_i(o)
\le
\varepsilon+\log_2\frac{1-p_{\max}}{1-2^\varepsilon p_{\max}}
\qquad\text{for every }i,o.
\]
Without the denominator condition, no finite prior-independent true-odds cap follows from a finite PML cap.
\end{proposition}

\begin{proof}
Bayes gives $\pi(i\mid o)=pm$.
Substitution into the definition of $R_i$ yields $m(1-p)/(1-pm)$.
For $m\ge1$, this expression is increasing in both $m$ and $p$ while $pm<1$, so the stated caps give the bound; for $m<1$ it is at most one and cannot enlarge the upper bound.
\end{proof}

\begin{proposition}[One bit of PML can coexist with infinite true-odds inflation]\label{prop:deterministic-reveal}
Let $I$ be uniform on two identities and let $O=I$ reveal the identity deterministically.
For the observed identity, $m_i(i)=2$, so $\ell(i)=1$ bit.
Its posterior is one, however, so its one-vs-complement posterior odds and $R_i(i)$ are infinite.
\end{proposition}

\begin{remark}[Why the terminology matters]
Pairwise inferential frameworks such as Pufferfish compare the posterior odds of one secret against another secret with their prior odds, equivalently a pairwise likelihood ratio.
The identity-vs-mixture factor $P_i/P$ is valuable and operationally meaningful, but it is not that pairwise odds ratio.
When the prior is tiny the two factors can be numerically close; the approximation is not a theorem and must not replace Proposition~\ref{prop:mass-to-odds}.\cite{KiferMachanavajjhala2014Pufferfish}
\end{remark}

\paragraph{Relation to PML and MaxL.}
The realized posterior-mass inflation $\ell(o)=\max_i\log_2(P_i(o)/P(o))$ is exactly PML at outcome $o$ under the declared prior.
Maximal leakage is the base-$2$ log-moment of $2^{\ell(O)}$ and yields the standard $\log_2(\alpha K)$ exact-identification sizing rule.
The endpoint-bridge note owns those algebraic identities and now also records the true-odds separation.\cite{series:endpointbridge,series:bossfightA}

\section{Closure properties}
\subsection{Post-processing}
\begin{theorem}[Post-processing]\label{thm:postprocessing}
Let $Z=g(O)$ be any (possibly randomized) function of $O$.
Then for every realized value $z$,
\[
\ell_Z(z)\le \operatorname*{ess\,sup}_{o:\,K(z\mid o)>0}\,\ell_O(o),
\]
where $K$ is the Markov kernel of $g$ (for deterministic $g$, this is $\sup_{o:\,g(o)=z}\ell_O(o)$).
In particular, the exact PML cap and the $(\varepsilon,\delta)_{\mathrm{HS}}$ envelope are preserved under post-processing.
\end{theorem}

\begin{proof}
Fix $z$ with $P_Z(z)>0$. For each $i$, write the Bayes factor $B_i(o)\defeq P_i(o)/P(o)$.
Then
\[
\frac{P_i(Z=z)}{P(Z=z)}=\frac{\sum_o P(o)B_i(o)\,K(z\mid o)}{\sum_o P(o)\,K(z\mid o)}=\EE_P\!\bigl[B_i(O)\mid Z=z\bigr],
\]
where $K$ is the Markov kernel of $g$ (randomized post-processing) and the conditional expectation is under the mixture $P$.
Therefore
\[
\max_i \frac{P_i(Z=z)}{P(Z=z)} \le \EE_P\!\Bigl[\max_i B_i(O)\,\Bigm|\, Z=z\Bigr].
\]
Since a conditional expectation is bounded by the essential supremum on the conditioning event,
\[
\EE_P\!\Bigl[\max_i B_i(O)\,\Bigm|\, Z=z\Bigr]
\le \operatorname*{ess\,sup}_{o:\,K(z\mid o)>0}\,\max_i B_i(o)
=2^{\operatorname*{ess\,sup}_{o:\,K(z\mid o)>0}\,\ell_O(o)}.
\]
Taking logs yields the stated bound.
For $(\varepsilon,\delta)$, note that \eqref{eq:hockeystick} is the $E_{2^\varepsilon}$ (hockey-stick) divergence bound $E_{2^\varepsilon}(P_i\|P)\le\delta$, and $E_\gamma$ satisfies data processing under Markov kernels, so the same $(\varepsilon,\delta)$ pair holds for $Z$ \cite{LiuCuffVerdu2017}.
\end{proof}

\subsection{Sequential composition under interaction}
Consider an interactive transcript $O_{1:T}$ where each step may depend on the past.
Write $\ell_t$ for the realized PML / posterior-mass inflation after observing prefix $O_{1:t}$.

\begin{theorem}[Sequential composition for $\varepsilon_\infty$]\label{thm:epsinf_comp}
If each step $t$ satisfies the conditional per-step bound
\[
\max_i \log_2\frac{P_i(O_t\mid O_{1:t-1})}{P(O_t\mid O_{1:t-1})}\ \le\ \varepsilon_t
\quad \text{almost surely},
\]
then the full transcript has PML / posterior-mass cap $\sum_{t=1}^T \varepsilon_t$.
\end{theorem}

\begin{proof}
Let $B_{i,t}\defeq \frac{P_i(O_{1:t})}{P(O_{1:t})}$ be the Bayes factor for identity $i$ after $t$ steps.
By the chain rule,
\[
B_{i,T}=\prod_{t=1}^T \frac{P_i(O_t\mid O_{1:t-1})}{P(O_t\mid O_{1:t-1})}.
\]
Taking $\log_2$, maximizing over $i$, and applying the per-step bound yields
$\ell(O_{1:T})=\log_2\max_i B_{i,T}\le \sum_{t=1}^T \varepsilon_t$ almost surely.
\end{proof}

\begin{theorem}[Sequential composition for $(\varepsilon,\delta)$ checkpoints]\label{thm:epsdelta_comp}
Consider an interactive transcript $O_{1:T}$.
Assume that for each step $t$ there exist numbers $(\varepsilon_t,\delta_t)$ such that, almost surely over the public history $O_{1:t-1}$,
for every event $A$ and every identity $i$,
\begin{equation}\label{eq:step_epsdelta}
P_i(O_t\in A\mid O_{1:t-1})\ \le\ 2^{\varepsilon_t}\,P(O_t\in A\mid O_{1:t-1})\ +\ \delta_t,
\end{equation}
where $P(\cdot\mid O_{1:t-1})$ is the mixture (posterior) baseline at step $t$.
Then the full transcript satisfies, for every event $B$ and every identity $i$,
\[
P_i(O_{1:T}\in B)\ \le\ 2^{\varepsilon}\,P(O_{1:T}\in B)\ +\ \delta_{\mathrm{tot}},
\qquad \varepsilon:=\sum_{t=1}^T \varepsilon_t,\quad
\delta_{\mathrm{tot}}:=\sum_{t=1}^T \delta_t\,2^{\sum_{s>t}\varepsilon_s}.
\]
In particular, if $\varepsilon_t\ge 0$ for all $t$, then $\delta_{\mathrm{tot}}\le 2^{\varepsilon}\sum_t\delta_t$.
\end{theorem}

\begin{proof}
We prove the bound by induction on $T$.
For $T=1$ it is \eqref{eq:step_epsdelta}.
Assume it holds up to $T-1$.
Fix an event $B$ on $O_{1:T}$ and write its sections
$B_h:=\{o_T:\ (h,o_T)\in B\}$ for histories $h=o_{1:T-1}$.
By the tower property,
\[
P_i(B)=\EE_{P_i}\bigl[P_i(O_T\in B_{O_{1:T-1}}\mid O_{1:T-1})\bigr].
\]
Applying \eqref{eq:step_epsdelta} with $A=B_{O_{1:T-1}}$ yields
\[
P_i(B)\le 2^{\varepsilon_T}\,\EE_{P_i}[\phi(O_{1:T-1})]+\delta_T,
\qquad \phi(h):=P(O_T\in B_h\mid O_{1:T-1}=h).
\]
To relate $\EE_{P_i}[\phi(O_{1:T-1})]$ to $\EE_P[\phi(O_{1:T-1})]=P(B)$, use the layer-cake identity
$\phi=\int_0^1 \one\{\phi\ge u\}\,du$ to reduce the claim to events and apply the induction hypothesis to each threshold set $\{\phi\ge u\}$.
By the induction hypothesis, with the same suffix-amplified slack rule applied to the first $T-1$ steps,
\[
\EE_{P_i}[\phi(O_{1:T-1})]\le 2^{\sum_{t<T}\varepsilon_t}\,\EE_P[\phi(O_{1:T-1})]+\sum_{t<T}\delta_t\,2^{\sum_{r=t+1}^{T-1}\varepsilon_r}.
\]
Multiplying this prefix slack by the final-step factor $2^{\varepsilon_T}$ and then adding $\delta_T$ gives
\[
\delta_T+\sum_{t<T}\delta_t\,2^{\varepsilon_T+\sum_{r=t+1}^{T-1}\varepsilon_r}
=\sum_{t=1}^T\delta_t\,2^{\sum_{s>t}\varepsilon_s},
\]
which is the stated $\delta_{\mathrm{tot}}$.
\end{proof}

\begin{corollary}[Homogeneous checkpoints]\label{cor:homog_checkpoints}
If the per-step bound \eqref{eq:step_epsdelta} holds with the same pair $(\varepsilon_t,\delta_t)\equiv(\varepsilon,\delta)$ for all $t$ (uniformly over histories), then for every event $B$ and every identity $i$,
\[
P_i(O_{1:T}\in B)\ \le\ 2^{T\varepsilon}\,P(O_{1:T}\in B)\ +\ \delta\,\sum_{j=0}^{T-1} 2^{j\varepsilon}
\ =\ 2^{T\varepsilon}P(B)+\begin{cases}
\delta\,\dfrac{2^{T\varepsilon}-1}{2^{\varepsilon}-1}, & \varepsilon\neq 0,\\
T\delta, & \varepsilon=0.
\end{cases}
\]
In particular, if $\varepsilon\ge 0$ then the slack is at most $2^{T\varepsilon}\,T\delta$.
\end{corollary}

\begin{remark}[Anytime-valid monitoring]
The per-identity Bayes factor process $B_{i,t}=2^{L_i(O_{1:t})}$ is a nonnegative martingale under the mixture law.
This connects posterior-mass inflation to \emph{e-processes} and time-uniform bounds (Ville-style inequalities), and to privacy odometers/filters in adaptive DP composition.
See \cite{HowardRamdas2018,RogersOdometers2016}.
\end{remark}

\section{Operational consequences}
Posterior-mass/PML caps imply concrete limits on what any adversary can do with the observation.

\begin{lemma}[Suspect sets must have prior mass]\label{lem:suspect}
For any threshold $\tau>0$ and any observation $o$,
\[
\pi\bigl(\{i:\ L_i(o)\ge \log_2\tau\}\bigr)\ \le\ \frac{1}{\tau}.
\]
In particular, if $\ell(o)\le \varepsilon$ then every identity $i$ satisfies $\pi(i\mid o)\le 2^\varepsilon \pi(i)$.
\end{lemma}

\begin{proposition}[Shortlists and credible sets]\label{prop:shortlist}
Suppose the mechanism has an $\varepsilon_\infty$ posterior-mass/PML cap.
Let $S(o)\subseteq\{1,\dots,N\}$ be any post-processed shortlist output on observation $o$.
If $S(o)$ is required to have posterior coverage $\pi(S(o)\mid o)\ge 1-\beta$ for all $o$, then $\pi(S(o))\ge 2^{-\varepsilon}(1-\beta)$ for all $o$.
Under a uniform prior this forces $|S(o)|\ge (1-\beta)2^{-\varepsilon}N$.
\end{proposition}

\begin{proof}
Fix any observation $o$ with $P(o)>0$.
By the $\varepsilon_\infty$ posterior-mass/PML cap, for every identity $i$ we have $\pi(i\mid o)\le 2^{\varepsilon}\pi(i)$.
Summing over $i\in S(o)$ gives
\[
\pi(S(o)\mid o)=\sum_{i\in S(o)}\pi(i\mid o)\ \le\ 2^{\varepsilon}\sum_{i\in S(o)}\pi(i)=2^{\varepsilon}\pi(S(o)).
\]
If $\pi(S(o)\mid o)\ge 1-\beta$, then rearranging yields $\pi(S(o))\ge 2^{-\varepsilon}(1-\beta)$.
For the uniform prior, $\pi(S(o))=|S(o)|/N$, hence $|S(o)|\ge (1-\beta)2^{-\varepsilon}N$.
\end{proof}

\begin{corollary}[Exact identification is bounded]\label{cor:identification}
Assume a uniform prior on $\{1,\dots,N\}$ and an $\varepsilon_\infty$ posterior-mass/PML cap.
Then for any estimator $\widehat I=\widehat I(O)$,
\[
\PP(\widehat I=I)\ \le\ \frac{2^{\varepsilon}}{N}.
\]
\end{corollary}

\begin{proof}
Condition on $O=o$.
The best success probability is the MAP mass $\max_i\pi(i\mid o)$, which is at most $2^{\varepsilon}/N$ under the $\varepsilon_\infty$ posterior-mass/PML cap and the uniform prior.
Average over $o$.
\end{proof}


\begin{corollary}[Identification under $(\varepsilon,\delta)$]\label{cor:identification_epsdelta}
Assume a uniform prior on $\{1,\dots,N\}$ and an $(\varepsilon,\delta)_{\mathrm{HS}}$ envelope in the sense of \Cref{def:epsinf}.
Then for any estimator $\widehat I=\widehat I(O)$,
\[
\PP(\widehat I=I)\ \le\ \frac{2^{\varepsilon}}{N}+\delta.
\]
\end{corollary}

\begin{proof}
Let $A_i=\{o:\ \widehat I(o)=i\}$.
Under the uniform prior, $\PP(\widehat I=I)=\frac{1}{N}\sum_{i=1}^N P_i(A_i)$.
Apply \eqref{eq:hockeystick} to each $A_i$ to obtain $P_i(A_i)\le 2^\varepsilon P(A_i)+\delta$.
Averaging over $i$ and using $\sum_i P(A_i)=1$ yields the claim.
\end{proof}



\begin{proposition}[Property inference is also capped]\label{prop:property}
Assume an $\varepsilon_\infty$ posterior-mass/PML cap. For any subset $E\subseteq\{1,\dots,N\}$ and any observation $o$ with $P(o)>0$,
\[
\pi(E\mid o)\ \le\ 2^{\varepsilon}\,\pi(E).
\]
Equivalently, no post-processed binary property of the identity can have its posterior probability inflated by more than $2^{\varepsilon}$.
\end{proposition}

\begin{proof}
Sum the pointwise bounds $\pi(i\mid o)\le 2^{\varepsilon}\pi(i)$ over $i\in E$.
\end{proof}

\section{Practical conversions from common contracts}
This section records two conversion lemmas that let mechanism-specific analyses ``plug into'' posterior-mass inflation.
The TV-equalization form is the same interface used by our TV-privacy padding compilers and dual-certificate tooling \cite{series:knapsacktransport}.

\subsection{TV equalization to a smoothed reference}
Let $Q$ be a public reference distribution on outcomes with $Q_{\min}:=\min_o Q(o)>0$.
In many anonymity channels, a mechanism exports a \emph{common-reference} guarantee: both the per-identity output law $P_i$ and the public marginal $P=\sum_i \pi(i)P_i$ are close to the same $Q$ (e.g.\ via equalization, smoothing, or explicit cover traffic).
This suffices to bound the identity-vs-mixture Bayes factors $P_i/P$ that define posterior-mass inflation.

\begin{lemma}[Common-reference TV-to-likelihood-ratio]\label{lem:tv_to_lr}
Let $P$ and $R$ be distributions on a finite alphabet, and let $Q$ satisfy $Q_{\min}>0$.
If $\TV(P,Q)\le \delta$ and $\TV(R,Q)\le \delta$ with $\delta<Q_{\min}$, then for all outcomes $o$,
\[
\left|\log_2\frac{P(o)}{R(o)}\right|\ \le\ \log_2\frac{Q_{\min}+\delta}{Q_{\min}-\delta}.
\]
\end{lemma}

\begin{proof}
Because total variation is the supremum event-probability gap, the singleton event $\{o\}$ gives $|P(o)-Q(o)|\le\delta$, hence $P(o)\le Q(o)+\delta$.
Likewise $R(o)\ge Q(o)-\delta$.
Therefore
\[
\frac{P(o)}{R(o)}\ \le\ \frac{Q(o)+\delta}{Q(o)-\delta}
\ \le\ \frac{Q_{\min}+\delta}{Q_{\min}-\delta},
\]
where the last step uses that $(x+\delta)/(x-\delta)$ is decreasing for $x>\delta$.
The reverse ratio $R(o)/P(o)$ is handled symmetrically.
\end{proof}

\begin{corollary}[Equalization $\Rightarrow$ per-output posterior-mass/PML cap]\label{cor:eq_to_epsinf}
If a mechanism satisfies $\TV(P_i,Q)\le \delta$ for all identities $i$ (with the same $Q$ and $\delta$),
then $\TV(P,Q)\le \delta$ for the public marginal $P=\sum_i \pi(i)P_i$.
Consequently, if $\delta<Q_{\min}$ then for every realized outcome $o$,
\[
\ell(o)=\max_i \log_2\frac{P_i(o)}{P(o)}\ \le\ \log_2\frac{Q_{\min}+\delta}{Q_{\min}-\delta}.
\]
\end{corollary}

\begin{proof}
By convexity of TV, $\TV(P,Q)=\TV(\sum_i\pi(i)P_i,Q)\le \sum_i \pi(i)\TV(P_i,Q)\le \delta$.
Apply \Cref{lem:tv_to_lr} with $P=P_i$ and $R=P$ and maximize over $i$.
\end{proof}

\begin{remark}
This conversion is conservative but operationally important: without a positive background mass ($Q_{\min}=0$), worst-case posterior-mass inflation can be infinite.
For anonymous DHTs, this is the ``rare committee'' failure mode.
When $Q_{\min}$ is too small to enforce, use a trimmed-support tail guarantee instead (next lemma).
\end{remark}

\begin{lemma}[Common-reference trimmed-support envelope]\label{lem:tv_trimmed}
Let $\tau>0$ and define the ``well-supported'' set $S_\tau\defeq\{o:\ Q(o)\ge \tau\}$.
If $\TV(P,Q)\le \delta$ and $\TV(R,Q)\le \delta$ with $\delta<\tau$, then for all $o\in S_\tau$,
\[
\left|\log_2\frac{P(o)}{R(o)}\right|\ \le\ \log_2\frac{\tau+\delta}{\tau-\delta}.
\]
Moreover, $P(S_\tau^c)\le Q(S_\tau^c)+\delta$ and $R(S_\tau^c)\le Q(S_\tau^c)+\delta$.
Thus a TV equalization contract yields a finite loss bound on $S_\tau$ and pushes the remaining probability mass into explicit ``rare outcome'' tail terms.
\end{lemma}

\subsection{Moment (R\'enyi) bounds to tail envelopes}
Mechanism analyses often bound conditional R\'enyi divergences or log-moments of the privacy loss.
Such bounds first give an actual loss-tail certificate by a one-line Chernoff argument; Proposition~\ref{prop:tail-vs-hs} then gives a hockey-stick envelope with the same slack.

\begin{lemma}[R\'enyi-to-tail conversion, bit units]\label{lem:renyi_tail}
Fix an identity $i$ and let $L_i(o)=\log_2\frac{P_i(o)}{P(o)}$ be the per-output privacy loss.
For $\alpha>1$, define the bit-valued R\'enyi divergence
\[
D_\alpha^{(2)}(P_i\|P)\ \defeq\ \frac{1}{\alpha-1}\log_2\sum_o P_i(o)^\alpha P(o)^{1-\alpha}.
\]
If $D_\alpha^{(2)}(P_i\|P)\le \rho_\alpha$, then for every bit threshold $\varepsilon$,
\[
P_i\!\bigl(L_i(O)\ge \varepsilon\bigr)\ \le\ 2^{(\alpha-1)(\rho_\alpha-\varepsilon)}.
\]
With $\delta=2^{(\alpha-1)(\rho_\alpha-\varepsilon)}$, this is a genuine tail certificate under $P_i$; it therefore also implies an $(\varepsilon,\delta)_{\mathrm{HS}}$ envelope by Proposition~\ref{prop:tail-vs-hs}. The reverse inference from an arbitrary hockey-stick row to this tail statement is invalid.
If an imported RDP accountant states a natural-log bound $D_\alpha^{(e)}(P_i\|P)\le r_\alpha$, use $\rho_\alpha=r_\alpha/\ln 2$ before applying this bit-unit tail formula.
\end{lemma}

\begin{proof}
By definition of the bit-valued R\'enyi divergence and the identity $2^{(\alpha-1)L_i}=\bigl(P_i/P\bigr)^{\alpha-1}$,
\[
\EE_{P_i}\!\left[2^{(\alpha-1)L_i(O)}\right]
=\sum_o P_i(o)\left(\frac{P_i(o)}{P(o)}\right)^{\alpha-1}
=\sum_o P_i(o)^{\alpha}\,P(o)^{1-\alpha}
=2^{(\alpha-1)D_\alpha^{(2)}(P_i\|P)}
\le 2^{(\alpha-1)\rho_\alpha}.
\]
Applying Markov's inequality under $P_i$ gives
\[
P_i(L_i\ge \varepsilon)=P_i\!\left(2^{(\alpha-1)L_i}\ge 2^{(\alpha-1)\varepsilon}\right)
\le 2^{(\alpha-1)(\rho_\alpha-\varepsilon)}.
\]
\end{proof}

\begin{remark}[Connection to RDP/moment accountants]
Lemma~\ref{lem:renyi_tail} is the same tail conversion that underlies R\'enyi differential privacy (RDP) and moment-accountant style envelopes, but applied to the identity-vs-mixture loss $L_i$ rather than neighboring datasets.
See \cite{MironovRDP2017} for the RDP framing, \cite{DworkRoth2014} for differential privacy background, and \cite{RogersOdometers2016} for adaptive composition tooling.
\end{remark}

\section{A worked audit-style example (telemetry to posterior-mass inflation)}\label{sec:worked_example}
Suppose a deployment fixes a discrete interface (output alphabet of size $K$) and a public background $Q$ with floor $Q_{\min}$ (or a trimmed floor $\tau$).
For each audited identity/class $i$ (or each equivalence class of identities), an auditor collects a window of $m$ i.i.d.\ samples from the interface and forms an empirical distribution $\widehat P_i$.
A conservative $L_1$ deviation bound of Weissman et al.\ yields a confidence radius for $\|\widehat P_i-P_i\|_1$ from finite samples \cite{Weissman2003L1}.
Translating to TV gives an upper confidence bound on $\TV(P_i,Q)$ when $Q$ is known.
Once all audited $P_i$ satisfy $\TV(P_i,Q)\le\delta$, \Cref{cor:eq_to_epsinf} converts this equalization certificate into a posterior-mass/PML cap against the public marginal.

\paragraph{Beyond TV radii (optional).}
If a deployment can support black-box hypothesis testing between $P_i$ and the mixture $P$, one can audit richer tradeoff curves (e.g., $f$-DP / tradeoff functions) rather than a single TV radius; see Askin et al.\ for recent black-box $f$-DP estimation and auditing methods \cite{AskinUSENIXSecurity2025}.

For a fully worked audit walk that threads telemetry, equalization, and posterior-mass-inflation semantics through the maintained example bundle, see Synthesis~17.\cite{series:workedexample} Exact theorem-root-to-review routing across that already-fixed endpoint card belongs to the maintained series-spine bridge (Synthesis~55) rather than to this note.\cite{series:seriesspines}

Release-facing packaging of the resulting claim surface belongs to Anonymity~C and Release~A rather than to this note.\cite{series:abomoinl,series:releaseA}

This ``telemetry $\rightarrow$ equalization $\rightarrow$ posterior-mass inflation'' pipeline is what Anonymity~B specializes to anonymous DHT committee traces, and what Anonymity~C records in an ABOM claim object \cite{series:profiling,series:abomoinl}.

\section{Release-facing unit and accountant validity boundary}\label{sec:unit_validity_boundary}
The public endpoint card exported by this paper is only valid after unit normalization has been checked.
A downstream anonymity card must name the transcript projection, prior/context declaration, budget form, and accountant units before it cites the composition or R\'enyi-tail rows above.
In particular, a bit-valued posterior-mass-inflation budget $\varepsilon$ cannot be combined with a natural-log RDP row, an $e^\varepsilon$ privacy-profile row, or a privacy-filter stopping rule until the row has been converted into the same base and bound type.
Adaptive or anytime monitors must also record whether their e-process / odometer bound is a pointwise $\varepsilon_\infty$ cap, a hockey-stick $(\varepsilon,\delta)$ checkpoint, or a moment row that still requires the conversion in Lemma~\ref{lem:renyi_tail}.

If any of these fields change---the exposure normal form, threat window, prior family, base of the accountant, R\'enyi order, stopping rule, trimming threshold, or rare-outcome tail policy---the old public endpoint card does not silently cover the new setting.
The correct maintenance action is a successor endpoint card or an explicitly signed conversion row, not a quiet reuse of the previous $\varepsilon$ symbol.

\section{How to consume this paper in the series}
Once a module exports either (i) an identity-vs-mixture likelihood-ratio cap, (ii) TV equalization to a smoothed reference, or (iii) a moment bound,
the posterior-mass/PML endpoint and its composition rules apply verbatim.
This allows each mechanism paper to avoid re-proving downstream operational consequences (shortlisting limits, identification bounds, etc.) and simply cite this paper.
For timing/termination-time contracts in anonymous overlay lookups, see the companion mechanism toolboxes \cite{series:schedcompiler,series:prefixcapacities}; for delegation-style overlays, see \cite{series:spectral}.


\paragraph{A minimal public posterior-mass / PML endpoint card.}
Table~\ref{tab:minimal-mass-card} is illustrative arithmetic, not evidence that any deployed channel attains the stated cap.

\begin{table}[t]
\centering
\small
\begin{tabular}{@{}p{0.33\linewidth}p{0.57\linewidth}@{}}
\toprule
\textbf{Public row} & \textbf{Pinned value}\\
\midrule
Declared observation surface & \texttt{exposure\_nf\_id = enf-committee-contact-bucket-v2}; \texttt{tw\_id = tw-24h-200-lookups} \\
Prior / context declaration & uniform over $N=4096$ eligible requesters in the declared shard/context \\
Endpoint contract form & $\varepsilon_\infty=0.75$ bits of PML / posterior-mass inflation on the declared surface \\
Pointwise posterior-mass cap & $\pi(i\mid o)\le2^{0.75}\pi(i)\approx1.681793\,\pi(i)$ for every $i,o$ \\
Guarded true-odds conversion & because $p_{\max}=1/4096$ and $2^{0.75}p_{\max}<1$, Proposition~\ref{prop:mass-to-odds} gives at most $0.750240221$ bits, a factor $1.682073$; this row is prior-dependent \\
Exact-identification cap & any estimator obeys $\Pr(\widehat I=I)\le2^{0.75}/4096\approx4.11\times10^{-4}$ \\
95\% shortlist floor & posterior coverage at least $0.95$ requires at least $0.95\cdot2^{-0.75}\cdot4096\approx2314$ identities \\
Evidence status & illustrative theorem drill only; no deployment channel, trace, or external-review certificate is supplied here \\
Downstream owner & Boss Fight~A / Anonymity~B / ABOM--OINL / Release~A for horizon accounting, anonymous-DHT specialization, and lineage packaging \\
\bottomrule
\end{tabular}
\caption{A minimal posterior-mass / PML endpoint card with a separately guarded true-odds conversion. The two quantities are not aliases.}
\label{tab:minimal-mass-card}
\end{table}

This card pins the observation surface, prior, endpoint representation, and valid corollaries without absorbing mechanism evidence, horizon accounting, or release objects. A public claim needs a source-bound channel witness and independent review in addition to this arithmetic.\cite{series:workedexample,series:seriesspines,series:bossfightA,series:profiling,series:abomoinl,series:releaseA}

\paragraph{One tiny worked-bundle endpoint drill.}
In the maintained worked bundle, the endpoint-facing line item \texttt{primary\_path\_declaration} in \texttt{example\_receipt.json} already fixes an exact root packet: \texttt{exposure\_nf\_id=sha256:5488040fe2921cbcda61a6d8f9a3887421abd33096e13c5dd1782d4436cbea8b}, \texttt{tw\_id=sha256:c1ff81fc84049a2160a21487f61a4565e03c5263f1437649c4299d1354b35c18}, witness tier \texttt{split\_identity\_destination\_surface}, qualifier \texttt{sep.nr1}, and budget value \texttt{0.0} bits per lookup for usage \texttt{Q<=500/24h}. In \texttt{example\_line\_item\_owner\_map.json}, that same endpoint root is carried by the public objects \texttt{[example\_primary\_surface\_manifest.json, example\_state\_decl.json, example\_replay\_plans.json]} with replay plan \texttt{eval-primary-v1}; \texttt{example\_abom.json} later binds the same primary-surface qualifier as \texttt{value=0.0} under separation class \texttt{sep.nr1}. That split is honest because this note owns the endpoint meaning of the declared zero-bit primary-path claim once the surface/context are fixed, while the maintained worked packet owns the concrete carrier row and Anonymity~C / Release~A own the later claim-object and signed-public wrappers.\cite{series:workedexample,series:seriesspines,series:abomoinl,series:releaseA}

\paragraph{One tiny endpoint drill.}
With the card values above,
\[
\Pr(\widehat I = I) \le \frac{2^{0.75}}{4096} \approx 4.11\times 10^{-4},
\]
by Corollary~\ref{cor:identification}. Likewise, Proposition~\ref{prop:shortlist} gives
\[
|S(o)| \ge 0.95\cdot 2^{-0.75}\cdot 4096 \approx 2314
\]
for every shortlist whose posterior coverage is at least $95\%$. A later terse paper can quote the card and this arithmetic without reopening the conversion lemmas or the release stack.

\paragraph{Same-card / refreshed-wrapper cutover.}
Suppose Release~$R$ and Release~$R'$ keep the same endpoint-card rows but refresh only the enclosing ABOM digest, signed receipt, or packet ordering because the source-only claim packet was reorganized. A later terse paper may then say \emph{``same endpoint card; release wrapper refreshed''} and stop at Anonymity~C / Release~A for lineage and signed-public-surface questions. If the declared surface, prior/context row, or budget form changes, reopen this note instead.\cite{series:abomoinl,series:releaseA}

\paragraph{One tiny wrapper-refresh drill.}
Keep the endpoint card in Table~\ref{tab:minimal-mass-card} fixed, but replace only the surrounding ABOM digest / signed-release receipt because the public packet now carries a clearer omission note. Then a later terse paper may honestly reuse the same arithmetic above and say only that the carrier was refreshed. The identification cap and shortlist floor stay owned by the same public object; only the enclosing release wrapper changed.

\paragraph{One tiny same-number cutover drill.}
Keep the same published budget number $\varepsilon_\infty=0.75$ bits, but quietly change the prior/context row from ``uniform over $N=4096$ eligible requesters in the declared shard'' to ``uniform over $N=1024$ eligible requesters in a narrower shard.'' Then the old endpoint card no longer licenses the same operational arithmetic, even though the budget number itself has not moved: the exact-identification cap becomes
\[
\frac{2^{0.75}}{1024} \approx 1.64\times 10^{-3},
\]
which is four times larger than the old $4.11\times 10^{-4}$ cap, and the $95\%$ shortlist floor drops to
\[
0.95\cdot 2^{-0.75}\cdot 1024 \approx 579.
\]
The same fail-closed rule applies if a later note reuses ``$0.75$ bits'' while changing ENF-ID or TW-ID: the budget number by itself is not the public claim object. A later terse paper may reuse the old arithmetic only while the declared observation surface, threat window, prior/context declaration, and endpoint budget form all stay fixed.

\paragraph{Minimal reopen ladder.}
Once the live question is no longer ``what exact endpoint anonymity card was declared here?'', the first reopen owner should be named explicitly rather than inferred. Reopen Boss Fight~A when the repeated-use horizon, mixed witness stack, or spending rule changes while the same endpoint semantics stay in force. Reopen Anonymity~B when the anonymous-DHT equalization witness, secret scope, or committee-trace specialization changes under the same endpoint metric. Reopen Anonymity~D when the live issue is the delegation witness tuple or witness-to-budget translation that feeds the already-declared anonymous-DHT card. Reopen State~1 when state-surface drift changes which prior/context or exposed control-plane surface the endpoint card is about. Reopen Operational~A or Operational~B when retry-discipline, replay-plan, or runtime-conformance questions change the declared projection itself. Reopen Eval~1 when the missing issue is notarized evaluator evidence or anti-grinding attempt logging. Reopen Anonymity~C and Release~A only when the live question has widened to claim-object packaging, signed public receipts, or cross-release lineage rather than endpoint semantics. Reopen Synthesis~17 when the exact downstream public/on-request packet for the endpoint card must be restated, and reopen Synthesis~55 when theorem-root-to-review routing across that already-fixed card is the missing object. This note therefore exports one compact stop rule: quote the endpoint card only while the declared surface, threat window, prior/context, and budget representation are the same public object.\cite{series:workedexample,series:seriesspines,series:bossfightA,series:profiling,series:spectral2budget,series:state1,series:operationalA,series:operationalB,series:eval1,series:abomoinl,series:releaseA}
\section{Takeaways}
\begin{itemize}[leftmargin=*]
\item \textbf{Endpoint metric:} posterior-mass inflation is a posterior-probability factor (PML), not a posterior-odds factor. A bounded claim must name the transcript projection $X$ and declared prior/context. It is a suitable endpoint only when the operational question and observer surface are stated.
\item \textbf{Composable by default:} sequential use over epochs composes additively in log posterior-mass factors (and via standard odometer/filter reasoning when stopping rules are adaptive). Mechanism papers in this archive should export per-epoch or per-step budgets and cite this paper for the endpoint/composition rule, while Boss Fight~A owns the repeated-use horizon interpretation.\cite{series:bossfightA}
\item \textbf{Operational conversions:} PML / posterior-mass caps induce (i) shortlist limits and (ii) identification sizing rules of the form $\log_2(\alpha K)$, enabling ``privacy-to-knobs'' compilation without re-deriving decision-theoretic consequences (see also Synthesis~5).\cite{series:endpointbridge}
\item \textbf{Claim surface, not release packaging:} the key inputs are the transcript projection, prior/context declaration, and budget representation. ABOM/OINL (Anonymity~C) is the canonical carrier for committing to those semantics in a source-only release; Release~A then packages lineage, anti-equivocation, and signed public release objects.\cite{series:abomoinl,series:releaseA}
\end{itemize}

\begin{thebibliography}{99}\small

\bibitem{series:contractobjects}
Anonymous.
\newblock Contract Objects for Anonymous-DHT Privacy Budgets and Claims.
\newblock Series synthesis note (Synthesis~0), 2026. (This archive.)


\bibitem{series:traceifaces}
Anonymous.
\newblock Trace Interfaces for Anonymous DHTs: Observation Tiers, Committee-Contact Traces, and Exposure Normal Forms.
\newblock Series synthesis note (Synthesis~3), 2026. (This archive.)

\bibitem{series:threatwindows}
Anonymous.
\newblock Threat Models and Repetition Windows for Anonymous DHTs: A Short Declaration Checklist for Referee-Grade Budget Claims.
\newblock Series synthesis note (Synthesis~4), 2026. (This archive.)

\bibitem{series:endpointbridge}
Anonymous.
\newblock Endpoint Metrics Bridge for Anonymous Overlays: Posterior-Mass Inflation, Pointwise Maximal Leakage, Maximal Leakage, and True Odds.
\newblock Series synthesis note, 2026. (This archive.)

\bibitem{SerjantovDanezis2002}
A.~Serjantov and G.~Danezis.
\newblock Towards an information theoretic metric for anonymity.
\newblock In \emph{Privacy Enhancing Technologies (PET)}, 2002.

\bibitem{RogersOdometers2016}
R.~Rogers, A.~Roth, J.~Ullman, and S.~Vadhan.
\newblock Privacy odometers and filters: pay-as-you-go composition.
\newblock In \emph{NeurIPS}, 2016. arXiv:1605.08294.

\bibitem{HowardRamdas2018}
S.~R.~Howard, A.~Ramdas, J.~McAuliffe, and J.~Sekhon.
\newblock Time-uniform Chernoff bounds via nonnegative supermartingales.
\newblock \emph{Probability Surveys}, 18:29--86, 2021. arXiv:1808.03204.

\bibitem{ZhaoSaeidianOechtering2026ContextAware}
H.~Zhao, S.~Saeidian, and T.~J.~Oechtering.
\newblock Context-aware privacy bounds for linear queries.
\newblock arXiv:2601.02855 [cs.IT], Jan.\ 2026. \url{https://arxiv.org/abs/2601.02855}.

\bibitem{AskinUSENIXSecurity2025}
\"{O}.~Askin, H.~Dette, M.~Dunsche, T.~Kutta, Y.~Lu, Y.~Wei, and V.~Zikas.
\newblock General-purpose $f$-DP estimation and auditing in a black-box setting.
\newblock In \emph{USENIX Security}, 2025. Preprint: arXiv:2502.07066. \url{https://arxiv.org/abs/2502.07066}.

% Common bibliography items for already-published papers in the anonymity series.
% Keep these stable so other papers can cite them without copying bib blocks.
% These refer to the companion archive already-published.zip.

\bibitem{series:schedcompiler}
Scheduling and Compiler Techniques for Metadata-Hiding Overlay Lookups.
\newblock Anonymity series paper (published; companion archive \texttt{already-published.zip}), 2026.
\newblock Wiki: \texttt{[[2026.01.22 - Mathematics: Scheduling and Compiler Techniques for Metadata-Hiding Overlay Lookups]]}.

\bibitem{series:prefixcapacities}
Prefix Capacities and Separation Cuts for Shared-Kernel Termination-Time Hiding.
\newblock Anonymity series paper (published; companion archive \texttt{already-published.zip}), 2026.
\newblock Wiki: \texttt{[[2026.01.25 - Mathematics: Prefix Capacities and Separation Cuts for Shared-Kernel Termination-Time Hiding]]}.

\bibitem{series:spectral}
Spectral Anonymity and Optimal Distinguishing Bounds for Random-Walk Delegation in (Possibly Directed) Overlays.
\newblock Anonymity series paper (published; companion archive \texttt{already-published.zip}), 2026.
\newblock Wiki: \texttt{[[2026.01.23 - Mathematics: Spectral Anonymity and Optimal Distinguishing Bounds for Random-Walk Delegation in (Possibly Directed) Overlays]]}.

\bibitem{series:stoptimeaddendum}
Stop-Time Padding Addendum (unique-only): Max-Mixture Separations and Tail-Sign Deadline Normal Forms.
\newblock Anonymity series paper (published; companion archive \texttt{already-published.zip}), 2026.
\newblock Wiki: \texttt{[[2026.01.26 - Mathematics: Stop-Time Padding Addendum, Max-Mixture Separations and Tail-Sign Deadline Normal Forms]]}.

\bibitem{SeriesManyTesting}
Many-Testing Lower Bounds for Low-Latency Privacy on Hierarchical Overlays.
\newblock Anonymity series paper (published; companion archive \texttt{already-published.zip}), 2026.

\bibitem{series:knapsacktransport}
Optimal Poset-Feasible Padding: Knapsack, Transport, and Dual Certificates for TV Privacy.
\newblock Anonymity series paper (published; companion archive \texttt{already-published.zip}), 2026.
\newblock Wiki: \texttt{[[2026.01.29 - Mathematics: Optimal Poset-Feasible Padding: Knapsack, Transport, and Dual Certificates for TV Privacy]]}.
% Common bibliography items for draft papers in the anonymity series.
% These are used for cross-references between the split papers.

\bibitem{series:profiling}
Anonymous DHT Profiling and the Equalization Contract.
\newblock Anonymity series Paper~B, 2026. (This archive.)

\bibitem{series:abomoinl}
ABOM and OINL for Anonymity Claims: Proof-Carrying Posterior-Mass Inflation Budgets Without Publishing Full Artifacts.
\newblock Anonymity series Paper~C, 2026. (This archive.)

\bibitem{series:releaseA}
Privacy as a Release Property: Proof-Carrying Anonymity Receipts for Anonymous DHTs.
\newblock Release and Destination series Paper~A, 2026. (This archive.)


\bibitem{MironovRDP2017}
I.~Mironov.
\newblock R\'enyi differential privacy.
\newblock In \emph{IEEE Computer Security Foundations Symposium (CSF)}, 2017. arXiv:1702.07476.

\bibitem{DworkRoth2014}
C.~Dwork and A.~Roth.
\newblock The algorithmic foundations of differential privacy.
\newblock \emph{Foundations and Trends in Theoretical Computer Science}, 9(3--4), 2014.

\bibitem{KiferMachanavajjhala2014Pufferfish}
D.~Kifer and A.~Machanavajjhala.
\newblock A framework for mathematical privacy definitions.
\newblock \emph{ACM Transactions on Database Systems}, 39(1):3:1--3:36, Feb.\ 2014.
\newblock doi:10.1145/2514681. (PDF: \url{https://users.cs.duke.edu/~ashwin/pubs/pufferfish_TODS.pdf}).

\bibitem{LiuCuffVerdu2017}
J.~Liu, P.~Cuff, and S.~Verd\'u.
\newblock $E_\gamma$-resolvability.
\newblock \emph{IEEE Transactions on Information Theory}, 63(5):2629--2658, May 2017.
\newblock doi:10.1109/TIT.2016.2642111. Preprint: arXiv:1511.07829.

\bibitem{NuradhaMeasuredHockeyStick2025}
T.~Nuradha, V.~Singh, and M.~M.~Wilde.
\newblock Measured hockey-stick divergence and its applications to quantum pufferfish privacy.
\newblock In \emph{Proc. IEEE International Symposium on Information Theory (ISIT)}, 2025.
\newblock doi:10.1109/ISIT63088.2025.11195501. arXiv:2501.12359. doi:10.48550/arXiv.2501.12359.
\newblock \url{https://arxiv.org/abs/2501.12359}.

\bibitem{Weissman2003L1}
T.~Weissman, E.~Ordentlich, G.~Seroussi, S.~Verd\'u, and M.~J.~Weinberger.
\newblock Inequalities for the $L_1$ deviation of the empirical distribution.
\newblock Hewlett-Packard Labs Technical Report HPL-2003-97R1, 2003.
\newblock (PDF: \url{https://www.hpl.hp.com/techreports/2003/HPL-2003-97R1.pdf}; mirror: \url{https://shiftleft.com/mirrors/www.hpl.hp.com/techreports/2003/HPL-2003-97R1.pdf}).




\bibitem{series:state1}
Anonymous.
\newblock State-Dependent Anonymity in Anonymous DHTs.
\newblock AnonDHT state series Paper~1, 2026. (This archive.)

\bibitem{series:operationalA}
Anonymous.
\newblock Schedule-Only Retries in Anonymous DHTs: SRP and Retry-Trace Equalization.
\newblock Operational series Paper~A, 2026. (This archive.)

\bibitem{series:operationalB}
Anonymous.
\newblock Auditable Trace Privacy with Abort.
\newblock Operational series Paper~B, 2026. (This archive.)

\bibitem{series:eval1}
Anonymous.
\newblock Notarized Anonymity Certificates: Attempt Logs, Spending Discipline, and Verifier Rules.
\newblock Evaluation series Paper~1, 2026. (This archive.)

\bibitem{series:artifactpolicy}
Anonymous.
\newblock Artifact policy for source-complete theorem notes.
\newblock Series note, 2026.


\bibitem{series:notation}
Anonymous.
\newblock Notation and Minimal Model for the One-True-Archive.
\newblock Series synthesis note (Synthesis~8), 2026. (This archive.)

\bibitem{series:bossfightA}
Anonymous.
\newblock Boss Fight Budgets: Maximal Leakage as a Repeated-Use Anonymity Contract.
\newblock Boss Fight series Paper~A, 2026. (This archive.)

\bibitem{series:workedexample}
Anonymous.
\newblock Worked Example: Instantiating a Tiered Reader-Privacy Receipt.
\newblock Synthesis~17, The-One-True-Archive (2026).

\bibitem{series:seriesspines}
Anonymous.
\newblock Series Spines: Math-to-Review Handoffs for Anonymous-DHT Papers.
\newblock Synthesis~55, The-One-True-Archive (2026).



% Added in rev0806 citation-closure hygiene pass.
\bibitem{series:spectral2budget}
Anonymous.
\newblock Spectral Delegation for Anonymous DHT Lookups: From mixing certificates to auditable per-step budgets.
\newblock Anonymity series Paper~D, The-One-True-Archive (2026).

\end{thebibliography}

\end{document}
